% chapters/introduction/document_structure.tex
% Sección 1.3: Estructura del documento (~0.5 páginas)
% ============================================================================
%
% GUÍA: Breve descripción de qué contiene cada capítulo.
%       Usar referencias cruzadas con \cref{ch:...} para los capítulos.
%
% Ejemplo:
%   El resto de este documento se organiza de la siguiente manera:
%
%   El \cref{ch:state_of_art} revisa los sistemas de corrección de exámenes
%   existentes y las tecnologías empleadas en el desarrollo.
%
%   El \cref{ch:requirements} recoge el análisis de requisitos funcionales
%   y no funcionales del subsistema backend.
%
%   El \cref{ch:design} presenta la arquitectura del sistema, el modelo de
%   datos, el diseño de la \ac{api} y los subsistemas principales.
%
%   El \cref{ch:implementation} describe la metodología de desarrollo, las
%   iteraciones realizadas y los aspectos técnicos más relevantes de la
%   implementación.
%
%   El \cref{ch:testing} detalla la estrategia de pruebas y los resultados
%   obtenidos en las pruebas unitarias, de integración y de carga.
%
%   Finalmente, el \cref{ch:conclusions} expone las conclusiones del
%   trabajo y las posibles líneas de trabajo futuro.
%
%   Además, se incluyen tres apéndices: el \cref{ap:global_requirements}
%   con los requisitos globales del sistema (frontend, backend y comunes),
%   el \cref{ap:deployment} con el manual de instalación y despliegue, y
%   el \cref{ap:detailed_results} con los resultados detallados de las
%   pruebas.
% ============================================================================

% TODO: Redactar la estructura del documento.
